\section{Category theoretic development}
\applabel{cat dev}

The amount of details in the theory of indexed symmetric data types is
high, but a few categorical abstractions keep the development
relatively clean. We start by recasting the parameterised initiality
theorems in this setting. Much of this is known but appears in
different places in the work of Fiore et
al.~\cite{fiore-plotkin-turi:ast}, and especially in the more recent
work with
Szamozvancev~\cite{fiore-szamozvancev:formal-soas,fiore-szamozvancev:soas-theory,szamozvancev:thesis}. We
provide references where available.
\begin{lemma}
  Let $P, F : \Cat C \to \Cat C$ be two endofunctors, and $\strength :
  PF \to FP$ a natural transformation. Then $F$ lifts along the
  forgetful functor $\fun{\pair xa}x$ to an endofunctor over
  $P$-algebras:
  \begin{gather*}
  \fun{\pair x{a : Px \to x}} \pair{Fx}{a_F : PFx \xto\strength FPx
    \xto{Fa} Fx}
  \\
  \fun{\parent{\pair xa \xto h \pair yb}}
  \parent{\pair{Fx}{a_F} \xto{Fh} \pair{Fy}{b_F}}
  \end{gather*}
\end{lemma}
The proof is by direct calculation, see Prop.~3.1.2 in
\citeauthor{szamozvancev:thesis}'s thesis. This type of result is
known from the early days of the concept, as in
\citeauthor{beck:distributive-laws}'s
work~\cite{beck:distributive-laws}.

When $P$ is a left adjoint $P \leftadjointto E$, we can swap
$P$-strengths for an endofunctor $F$ with $E$-costrengths for the same
endofunctor\footnote{We learned this fact from Marcelo
P.~Fiore in private communication.}:
\begin{lemma}[strength-costrength]
  Consider endofunctors $P, E, F : \Cat C \to \Cat C$. Let $P
  \leftadjointto E$ with unit $\eta : \Id \to EP$ and counit $\epsilon :
  PE \to \Id$. TFAE:
  \begin{itemize}
  \item $P$-strengths for $F$, i.e., natural transformations
    $\strength : PF \to FP$.
  \item $E$-costrengths for $F$, i.e., natural transformations
    $\costrength : FE \to EF$.
  \end{itemize}
  The equivalence is given by the following relationships:
  \begin{center}
    \diagram*{31} \qquad
    \diagram*{32}
  \end{center}
  and moreover satisfy:
  \begin{center}
    \diagram*{33}
    \qquad
    \diagram*{34}
  \end{center}
\end{lemma}
\begin{proof}
  By duality, we only need to show:
  \begin{itemize}
  \item That $\costrength$ is natural.
  \item That the strength associated with $\costrength$ is $\strength$.
  \item The diagram involving the units only ($\eta$ and $F\eta$).
  \end{itemize}

  The $\costrength$ is natural: it is a
  composition of natural transformations.

  We substitute the definition of $\costrength$ into the definition of
  $\strength$ and show it recovers $\strength$:
  \diagram{36}

  The diagram involving units only amounts to straightforward
  calculation:
  \diagram{35}
\end{proof}

We now prove a general version of parameterised initiality, already
prevalent in Fiore et al.'s work on abstract
syntax~\cite{fiore-szamozvancev:formal-soas,%
  fiore-szamozvancev:soas-theory,%
  fiore-ranchod:substructural-ast}, originally proposed by
\citeauthor{bird-paterson:generalised-folds-for-nested-datatypes}~\cite{%
  bird-paterson:generalised-folds-for-nested-datatypes}.

\begin{definition}
  Let $P, F : \Cat C \to \Cat C$ be endofunctors. An
  initial $F$-algebra $\roll : F\fix F \to \fix F$ is
  \emph{$P$-initial} when every $F$-algebra $f : Fx \to x$ has a
  unique morphism $\pfold f : P\fix F \to x$ satisfying: \diagram{37}
\end{definition}

\begin{theorem}[generalised parameterised initiality]%
  \theoremlabel{generalised parameterised initiality}
  Let $P \leftadjointto (E : \Cat C \to \Cat C)$ be an adjunction over
  a category $\Cat C$, with unit $\eta : \Id \to EP$ and counit
  $\epsilon : PE \to \Id$. Every initial algebra $\roll : F\fix F \to
  \fix F$, for an endofunctor $F : \Cat C \to \Cat C$ equipped with a
  natural transformation $\strength : PF \to FP$, is $P$-initial.
\end{theorem}

For completeness, we include the proof of this result in this
formulation.

\begin{proof}
  Let $f : Fx \to x$ be any $F$-algebra.  We equip $Ex$ with an
  $F$-algebra structure. It is more direct to work with the costrength:
  \[
  g : FEx \xto\costrength EFx \xto{Ef} Ex
  \]
  and so, by initiality, we have a unique morphism $h \definedby \fold
  g : \fix F \to Ex$ satisfying:
  \diagram{38}
  Now take the mate of this morphism:
  \[
  \pfold f :P\fix F \xto{Ph} PEx \xto{\epsilon} x
  \]
  This morphism satisfies the required property:
  \diagram{39}
  Now consider any morphism $k : P\fix F \to x$ making the following
  diagram commute:
  \diagram{40}
  We will show that its mate must be $h$, i.e.: $h : \fix F \xto{\eta}
  EP\fix F \xto{Ek} Ex$. Calculate:
  \diagram{41}
  By universality, $h$ is the mate of $k$, and so $k = \pfold f$.
\end{proof}

Therefore, if an endofunctor $F$ over $\carrier\Ctx$-indexed families
has a $(\gast-)$-strength, then its initial algebras are both
$(\gast-)$-initial and $(\gast\gast-)$-initial.

\begin{definition}
  Let $P = \triple{P : \Cat C \to \Cat C}{\eta}{\mu}$ be a monad, and $F :
  \Cat C \to \Cat C$ an endofunctor over $\Cat C$. A \emph{$P$-strength for $F$}
  is a natural transformation $\strength : PF \to FP$ satisfying:
  \begin{center}
    \diagram*{42}\qquad
    \diagram*{43}
  \end{center}
\end{definition}
This notion is sometimes called: a \emph{distributive law} between a
monad and a functor; a \emph{module}; or a \emph{monad endomorphism}.

The following notion generalises equivariant $F$-algebras as algebras
of the lifting of $F$ to $P$-algebras:
\begin{definition}
  Let $P = \triple{P : \Cat C \to \Cat C}{\eta}{\mu}$ be a monad over
  $\Cat C$, and $F : \Cat C \to \Cat C$ an endofunctor equipped with a
  $P$-strength $\strength P : PF \to FP$. A \emph{$P$-linear $F$-algebra}
  $\triple{x}{f}{a}$ consists of:
  \begin{itemize}
  \item an $F$-algebra $f : Fx \to x$; and
  \item a $P$-algebra $a : Pa \to a$;
  \end{itemize}
  such that $f$ is $P$-linear $f : \pair{Fx}{a_F} \to \pair{x}{a}$. A
  \emph{homomorphism} $h : \triple xfa \to \triple ygb$ of $P$-linear
  $F$-algebras is a morphism $h : x \to y$ that is simultaneously an
  $F$-algebra homomorphism and a $P$-algebra homomorphism.
\end{definition}

Since a $(\gast-)$-algebra is precisely a $\Ctx$-indexed $\G$-action
and a $(\gast-)$-linear map between them is precisely an equivariant
function, we have that the $(\gast-)$-linear $F$-algebras are
precisely the equivariant $F$-algebras. The capstone of this
development is the following theorem:

\begin{theorem}[general representation theorem]%
  \theoremlabel{general representation}
  Let $P = \triple{P : \Cat C \to \Cat C}{\eta}{\mu}$ be a monad over
  $\Cat C$, and $F : \Cat C \to \Cat C$ an endofunctor equipped with a
  $P$-strength $\strength : PF \to FP$. If an initial $F$-algebra
  $\roll : F\fix F \to \fix F$ is both $P$-initial and $P^2$-initial,
  then it is also the initial $P$-linear $F$-algebra. Its $P$-algebra
  structure is given by $P$-initiality as
  \(
  \alg \definedby \pfold\roll : P\fix F \to \fix F
  \)
  and the unique $F$-homomorphism $\fold f : \fix F \to x$ into any
  other $P$-linear $F$-algebra $\triple xfa$ is $P$-linear.
\end{theorem}

We chose this formulation as its proof follows our proof in the
simply-typed case more closely, which itself follows the argument in
abstract syntax with binding and
substitution~\cite{fiore-plotkin-turi:ast}. Our situation is
semantically well-behaved, and there are alternative theorems we could
use. For example, since $\Ctx$-indexed $\G$-actions are $E \definedby
(\lollybox-)$-algebras, we could use
\citeauthor{szamozvancev:thesis}'s Prop.~3.3.1
instead~\cite{fiore-szamozvancev:soas-theory}. Whichever formulation
we use, it is clear it proves the Indexed Representation
\theoremref{indexed representation}.

\begin{proof}
  We show that:
  \begin{itemize}
  \item $\alg : P\fix F \to \fix F$ is a $P$-algebra;
  \item $\roll$ is $P$-linear by the defining property of $P$-initiality;
  \item each $\fold f$ is $P$-linear.
  \end{itemize}

  We show the $P$-algebra unit axiom:
  \diagram{44}
  by recourse to mere initiality:
  \diagram{45}
  We show the $P$-algebra multiplication axiom:
  \diagram{47}
  by recourse to $P^2$ initiality:
  \diagram{46}
  Finally, take any $P$-linear $F$-algebra $\triple xfa$. By
  initiality there is at most one $P$-linear $F$-homomorphism, and
  exactly one $F$-homomorphism, into it. We show:
  \diagram{48}
  by recourse to $P$-initiality:
  \diagram{49}
  concluding the proof.
\end{proof}
